The morphism $q$ is projective and of finite presentation.  Apply generic
flatness to $q$ and to $\cA_i\to\cM$.  Geometric reducedness of the fibres of
$q$ and geometric integrality of the fibres of each $\cA_i$ then hold after
restricting to a nonempty open of $\cM$; this is the standard spreading of
these properties in a flat family \cite[\href{https://stacks.math.columbia.edu/tag/052A}{Tag 052A};
\href{https://stacks.math.columbia.edu/tag/0578}{Tag 0578};
\href{https://stacks.math.columbia.edu/tag/0559}{Tag 0559}]{stacks-project}.  This proves the
flatness and fibre assertions in (i), as well as the first assertion of (ii).

Since $q$ is projective and the $\cA_i$ are flat, their closed immersions
define two sections of the relative Hilbert scheme
$\operatorname{Hilb}_{\cI/\cM}$.  Its diagonal is closed.  The equalizer of
the two sections is therefore closed and omits the generic point.  Removing
it makes $(\cA_1)_x$ and $(\cA_2)_x$ distinct for every geometric $x$.

Give $\cA=\cA_1\cup\cA_2$ its scheme-theoretic union structure.  Let
$\mathcal J$ be the ideal of $\cA\cap\cI^{\mathrm{ne}}$ in
$\cI^{\mathrm{ne}}$.  The geometric generic fibre is reduced and consists
exactly of the two generic components, so $\mathcal J$ vanishes there.  The
image of $\operatorname{Supp}(\mathcal J)$ in the noetherian integral scheme
$\cM$ is constructible by Chevalley's theorem and does not contain its generic
point.  Its closure is therefore proper.  Remove this closure.  Then
$\mathcal J=0$, and
\[
 \cI^{\mathrm{ne}}=\cA\cap\cI^{\mathrm{ne}}
\]
over the remaining open.  The fibrewise reducedness, the distinctness and
geometric integrality of the two curve fibres, and this scheme-theoretic
equality prove (ii) and (iii).

Denote the remaining open by $\cM^\circ$.
The morphism $\cM^\circ\to\cB$ is open because it is the restriction of the
smooth morphism $\cM\to\cB$ \cite[\href{https://stacks.math.columbia.edu/tag/056G}{Tag 056G}]{stacks-project}.  Hence its image
$\cB^\circ$ is nonempty and open.  For $b\in\cB^\circ(\CC)$, the fibre
$U_b$ is a nonempty open of the integral surface $(X_b)_{\reg}$; it is
therefore dense, irreducible, and connected.  The fibrewise conclusions
(i)--(iii) give the final assertion.
\end{proof}

\section{Pair convolution detects \texorpdfstring{$\IC_X$}{IC(X)}}

Fix $b\in\cB^\circ(\CC)$, put $(P,X)=(P_b,X_b)$, and let $U=U_b$ as in
\cref{thm:uniform-pair-splitting}.
Let $\rho:S\to X$ be the minimal resolution and let $j:S\to P$ be its
composition with the inclusion.  Put
\[
 K=\IC_X,
 \qquad
 \widetilde K=Rj_*\QQ_S[2]=K\oplus\delta_0.
\]
Consider
\[
 \wt m:S\times S\longrightarrow P,
 \qquad (p,q)\longmapsto j(p)+j(q),
\]
and let $\sigma(p,q)=(q,p)$.  A superscript $-$ denotes the
anti-invariant direct summand under $\sigma$.

\begin{lemma}[Top cohomology on the split-pair locus]
\label[lemma]{lem:split-fibre-top}
After replacing $U$ by a dense open subset, put $U^+=-U$.  For $t=-x\in
U^+$, the fibre $F_t=\wt m^{-1}(t)$ has exactly four one-dimensional
irreducible components in its reduction: the strict transforms of
$(\cA_1)_x$ and $(\cA_2)_x$ and the two structural components
\[
 E_{\exc}\times\{s_t\},\qquad\{s_t\}\times E_{\exc}.
\]
Here $s_t$ is the unique point of $S$ above $t$.  The involution $\sigma$
exchanges the first two components and exchanges the last two.
Consequently
\[
 \dim\cH^{-2}(\widetilde K*\widetilde K)_{-,t}=2,
 \qquad
 \dim\cH^{-2}(K*K)_{-,t}=1.
\]
After a further shrinking of $U$, the latter spaces form the constant
rank-one local system on $U^+$.
\end{lemma}

\begin{proof}
Fix $t\in U^+$, write $x=-t\in U$, and let
$E_{\exc}=\rho^{-1}(0)$.  By definition, $C_x$ parametrizes ordered pairs
whose sum is $-x=t$.  Hence the fibre
$F_t=\wt m^{-1}(t)$ has precisely four irreducible curve components: the
strict transforms of $(\cA_1)_x,(\cA_2)_x$ and the two structural components
\[
 C_{1,t}=E_{\exc}\times\{s_t\},
 \qquad
 C_{2,t}=\{s_t\}\times E_{\exc},
\]
where $s_t$ is the unique point above $t$.  Indeed, away from
$E_{\exc}$ the resolution is an isomorphism, while a pair with one entry on
$E_{\exc}$ must have the other entry over $t$.  Thus there are no further
top-dimensional components.  The involution $\sigma$ exchanges
$(\cA_1)_x$ with $(\cA_2)_x$ and $C_{1,t}$ with $C_{2,t}$.

Proper base change identifies
\[
 \cH^{-2}(\widetilde K*\widetilde K)_t=H^2(F_t,\QQ).
\]
The top cohomology of a proper complex curve is freely generated by the
fundamental classes of the irreducible components of its reduction.  Thus
$H^2(F_t,\QQ)_-$ has the explicit basis
\[
 [\,(\cA_1)_x\,]-[\,(\cA_2)_x\,],
 \qquad [C_{1,t}]-[C_{2,t}].
\]
Here and below the same notation is used for a curve and its strict transform
in $F_t$.  Holomorphic transposition preserves complex orientation, so
\begin{equation}
 \dim\cH^{-2}(\widetilde K*\widetilde K)_{-,t}=2.
\label{eq:total-alt-rank}
\end{equation}

Expand the decomposition $\widetilde K=K\oplus\delta_0$:
\[
 \widetilde K*\widetilde K=(K*K)\oplus(K*\delta_0)\oplus(\delta_0*K)
      \oplus(\delta_0*\delta_0).
\]
At $t\ne0$, the last term has zero stalk.  The two middle terms are both
canonically $K$, and transposition exchanges them.  Since $t$ is a smooth
point of the surface $X$, their anti-invariant stalk in degree $-2$ is the
line generated by $[C_{1,t}]-[C_{2,t}]$.  Taking dimensions in this direct
sum decomposition gives
\begin{equation}
 \dim\cH^{-2}(K*K)_{-,t}=1.
\label{eq:principal-alt-rank}
\end{equation}
There is no Koszul sign: $K$ has geometric shift $[2]$, and
$(-1)^{2\cdot2}=1$.

After a topological stratification and a further shrinking of $U^+$, proper
stratified local triviality identifies the four component classes in nearby
fibres.  Hence the total anti-invariant local system is constant of rank two.
On $U^+$ the sheaf $K$ is the shifted constant sheaf $\QQ_{U^+}[2]$, so the
anti-invariant part of
$(K*\delta_0)\oplus(\delta_0*K)$ is itself a constant rank-one local system.
The direct-sum decomposition above therefore makes $\cH^{-2}(K*K)_-$ a
rank-one direct summand of a constant rank-two local system, with constant
complement.  Equivalently, its monodromy is trivial.  This proves the last
assertion without choosing a preferred lift of the class
$[\,(\cA_1)_x\,]-[\,(\cA_2)_x\,]$.
\end{proof}

\begin{lemma}[Vanishing away from the theta divisor]
\label[lemma]{lem:off-theta-vanishing}
There is a dense open subset $W\subset P\setminus X$ such that, for every
$t\in W$, the scheme $X\cap(t-X)$ is a smooth connected projective curve.
The involution $p\mapsto t-p$ acts as the identity on its top cohomology.  In
particular,
\[
 \cH^{-2}(K*K)_-|_W=0.
\]
\end{lemma}

